\documentclass[11pt,letterpaper]{article}

\usepackage[utf8]{inputenc}
\usepackage[T1]{fontenc}
\usepackage{lmodern}
\usepackage{geometry}
\usepackage{microtype}
\usepackage{booktabs}
\usepackage{longtable}
\usepackage{array}
\usepackage{enumitem}
\usepackage{amsmath}
\usepackage[round,authoryear]{natbib}
\usepackage{xcolor}
\usepackage{hyperref}

\geometry{margin=1.15in}
\hypersetup{colorlinks=true, linkcolor=black, citecolor=black, urlcolor=blue!55!black}

\setlist{itemsep=2pt, topsep=4pt}
\renewcommand{\arraystretch}{1.25}

\title{\vspace{-1.2cm}\textbf{Philosophy Engineering Is Not a Paradigm}\\[0.5em]
\large A test against Kuhn's criteria, and a present-tense answer to\\
Feyerabend's objection to Lakatos}

\author{Andrew H. Bond\\
\small Department of Computer Engineering, San Jos\'{e} State University\\
\small \texttt{andrew.bond@sjsu.edu}}

\date{\small September 2026 \\[0.4em]
\small\textit{Draft. Evidence class:} \texttt{exploratory}\textit{. See \S10.}}

\begin{document}
\maketitle

\begin{abstract}
\noindent
Philosophy Engineering has been described by its author as a new paradigm for
science in Kuhn's sense. This paper tests that description against Kuhn's own
criteria and rejects it. The programme supplies symbolic generalizations and a
changed account of what counts as a solved problem, and it has a small stock of
exemplars. It has no community, which for Kuhn is not a rhetorical shortfall but
the unit of analysis. It also possesses a specified procedure for translating a
legacy corpus into its own vocabulary, and the existence of that procedure is
evidence against incommensurability and therefore against revolution. The
anomaly the programme names, representation dependence, cannot produce a
Kuhnian crisis, because the reigning instruments score a representation-dependent
system as correct. Kuhn requires a paradigm to fail by its own standards, and
this one does not. What survives the negative result is smaller and better
supported. Lakatos distinguished a programme's hard core from its protective
belt and was criticized by Feyerabend for an appraisal available only in
hindsight. A claim ledger that fixes priority by repository ancestry and
requires dependencies to be declared before the test turns that partition into a
property of the record rather than a reconstruction, so the question of what a
refutation strikes acquires a present-tense answer. That is not a paradigm. It
is a change to the medium in which paradigms are recorded. The paper identifies
the inquiry arm positively, as a set of secondary rules in Hart's sense, and
states five predictions by which the account can be shown wrong.
\end{abstract}

\vspace{0.6em}
\noindent\textbf{Keywords:} Kuhn, Lakatos, Feyerabend, research programmes,
scientific infrastructure, preregistration, philosophy engineering, claim ledger

\section{Introduction}

Philosophy Engineering is the practice of taking a philosophical commitment and
building something that mechanically refuses to violate it. Its author has
described it as a new paradigm for science in the sense of Kuhn
\citep{kuhn1970}. This paper tests that description. The result is negative.

The test is worth running because the claim is cheap. Announcing a paradigm
shift costs nothing, cannot be refuted by any short argument, and is made
routinely by programmes that turn out to be ordinary. A discipline whose stated
rule is that a programme may not assert what its ledger cannot show is poorly
placed to make the cheapest claim in science and leave it unexamined. The
appropriate response is to state the criteria in advance, apply them, and report
what happens, including to the parts of the claim that fail.

Three results follow. First, the programme fails Kuhn's community criterion by a
wide margin, and it fails the incommensurability criterion in a way that counts
in its favour rather than against it, because the thing that fails the criterion
is a working translation procedure. Second, the anomaly on which the strongest
case rests cannot do the work Kuhn's account requires of anomalies, since a
system exhibiting it still passes the measurements the reigning approach uses.
Kuhn's crises begin when a paradigm fails by its own standards, and this
anomaly is invisible to those standards rather than damaging to them. Third, the
programme contains two arms with different verdicts, and the appearance of a
paradigm comes largely from running them together: the judgment arm has the
shape of a candidate theory change inside one subfield, while the inquiry arm is
not a theory at all and was never a candidate.

What remains after the negative result is a smaller claim with better support.
Lakatos separated a research programme's hard core from its protective belt and
argued that refutations ordinarily strike the belt \citep{lakatos1970}.
Feyerabend objected that the resulting appraisal was available only in hindsight,
which made it useless to a scientist deciding what to do now
\citep{feyerabend1975}. The objection has force because nothing in the ordinary
scientific record fixes which commitments were core before the test was run, so
the partition is reconstructed afterwards by someone with an interest in the
answer. A record that fixes priority by repository ancestry, requires
dependencies to be declared before the seal, and computes the consequences of a
refutation by traversing a graph, turns that partition into something the record
already contains. The question of what a refutation strikes then has an answer
in the present tense. This is a contribution to the appraisal of research
programmes, not a revolution, and the paper states five ways it could be shown
wrong.

\section{The claimant is badly placed to assess the claim}

Kuhn's own account makes a participant's declaration of a paradigm shift weak
evidence. Revolutions in \emph{Structure} are described as invisible from the
inside. The community's textbooks are rewritten so that the current position
appears to be the cumulative result of what came before, and the practitioners
who perform the rewriting do not experience themselves as doing so
\citep[ch.~11]{kuhn1970}. If revolutions are hard to see while standing in them,
the person best placed to see one is a historian at some distance, and the person
worst placed is the programme's founder.

There is a second difficulty. The word is used loosely. Masterman counted
twenty-one distinct senses of ``paradigm'' in the first edition of
\emph{Structure} \citep{masterman1970}, and Kuhn accepted enough of the criticism
to separate two of them in later work: the disciplinary matrix, meaning
everything a community shares, and the exemplar, meaning a concrete problem
solution that trains practitioners to recognize new problems as being of the same
kind \citep{kuhn1974}. A claim to be a paradigm is therefore ambiguous until the
sense is fixed. Both senses are tested below.

The procedural safeguard adopted here is the one the programme applies to
itself. Under PE-CLS-1.0 no document may cite a claim above its evidence class
\citep{pecls}. The thesis of this paper is supported by one programme, observed
by its own author, and has not been tried elsewhere. Its class is
\texttt{exploratory}, and no other document in the corpus may cite it as more
than that. Section~10 states what the record cannot currently show.

\section{What Kuhn requires}

Kuhn's criteria are stated here in the form in which they will be applied.

\textbf{The disciplinary matrix.} A paradigm in the broad sense is what the
members of a scientific community share. Kuhn decomposes it into symbolic
generalizations deployed without argument, models and metaphysical commitments
that supply preferred analogies, values that govern what counts as a good
solution, and exemplars \citep[Postscript]{kuhn1970}.

\textbf{Exemplars.} In the narrow sense a paradigm is a concrete problem
solution, transmitted through training, that lets a practitioner see an
unfamiliar situation as similar to a familiar one \citep{kuhn1974}. Kuhn treats
this as the primary sense. Exemplars are numerous, they are taught, and they do
their work by resemblance rather than by rule.

\textbf{The community as the unit.} Kuhn's definition is openly circular. A
paradigm is what a scientific community shares, and a scientific community
consists of those who share a paradigm \citep[Postscript]{kuhn1970}. The
circularity is not an embarrassment in his account; it fixes the unit of
analysis. Paradigms are properties of communities of practitioners. They are not
properties of positions, theories, or individuals.

\textbf{Normal science and anomaly.} Normal science is puzzle-solving conducted
inside the matrix, where the existence of a solution is assumed and the
practitioner's skill is at issue. An anomaly is a result the paradigm cannot
assimilate. Anomalies alone do not produce crises. A crisis arises when the
paradigm repeatedly fails at puzzles it has itself certified as solvable, and the
failure becomes visible to the community using the paradigm's own instruments
\citep[chs.~6--8]{kuhn1970}.

\textbf{Incommensurability.} Rival paradigms lack a neutral vocabulary in which
their claims can be compared without loss. Kuhn's later statements narrow this
to taxonomic incommensurability, a failure of intertranslatability among the
kind terms of the two theories. The relevant consequence for this paper is
negative and sharp: there is no algorithm that converts the claims of one
paradigm into the claims of another with their evidential standing intact. Where
such a procedure exists, the two bodies of work are not separated by a
revolution.

\section{The case for}

The case that Philosophy Engineering is a paradigm is stronger than its brevity
here suggests, and it should be put at full strength before it is tested.

The programme names an anomaly. Systems deployed in medicine, law, hiring, and
content moderation change their verdicts when the description of a case changes
in ways that preserve its content. Order of presentation, party names, active
against passive voice, choice of synonym, and units of measure all move outputs
that ought to be fixed by the situation rather than by its wording. The
Foundation document collects the failure modes and treats them as instances of a
single defect rather than as unrelated bugs \citep{foundation}.

The programme supplies symbolic generalizations. The Epistemic Invariance
Principle requires a judgment procedure to return equivalent outputs for
representations related by a declared transformation group. The Bond Invariance
Principle requires that transformations preserving the relational structure of a
case leave the evaluation unchanged, and yields an accountability theorem: any
two differing judgments differ because the relational structure differs, because
the evaluative parameters differ, or both, with no fourth option
\citep{foundation}. Both are stated formally and proved.

The programme changes what counts as a good solution. Under the reigning
approach a classifier is good when its aggregate accuracy is high. Under this
one a judgment procedure is good when it is invariant under the declared group
and still discriminates between cases that differ. Invariance alone is
satisfied perfectly by a constant function, so the second condition is not
decoration. A change in the criteria of adequacy is exactly what Kuhn lists among
the components of a disciplinary matrix.

The programme has exemplars. A failed invariance test produces a witness, a
minimal transformation that flips the verdict, and learning to read a witness
teaches a practitioner what kind of object a representation-dependence failure
is. The derivation of the Klein four-group from Hohfeld's jural relations
\citep{hohfeld1913} shows how a deontic vocabulary becomes a symmetry claim. The
withdrawal episode recorded in the corpus shows how a distinction between
unsupported and refuted operates on a dependency graph \citep{pecls}. Brownfield
conversion shows what to do with a corpus written before the method existed
\citep{pebrw}. These are concrete problem solutions of the kind Kuhn describes,
and a practitioner who has worked through them sees new cases differently.

Finally, the field the programme addresses talks like a field in crisis. Its
practitioners describe alignment as unsolved, its evaluation methods as
saturated, and its failures as poorly understood. Crisis rhetoric is not crisis,
but Kuhn treats the community's own sense of malaise as one of the marks.

\section{The test}

\begin{longtable}{p{0.26\textwidth} p{0.10\textwidth} p{0.55\textwidth}}
\toprule
\textbf{Kuhn's criterion} & \textbf{Verdict} & \textbf{Basis} \\
\midrule
\endhead
Symbolic generalizations & Partial & EIP and BIP are stated and proved. Kuhn's generalizations are deployed by a community without argument; these are argued for by one author. \\
Values and criteria of adequacy & Pass & Invariance with discrimination replaces aggregate accuracy. This is a change in what counts as a solved problem. \\
Exemplars & Partial & Five or six exist. Kuhn's paradigms supply hundreds, transmitted by training. There is no training pipeline. \\
Community & \textbf{Fail} & One author, one department, a small number of students, a chat server. The unit of analysis is absent. \\
Normal science & Weak & Puzzle-solving inside the matrix occurs, set and solved by the same person. \\
Crisis in a predecessor & \textbf{Fail} & The anomaly does not register as failure on the reigning instruments. See below. \\
Incommensurability & \textbf{Fail} & A specified translation procedure exists and has been run. See below. \\
\bottomrule
\end{longtable}

\textbf{Community.} This is the decisive failure and it does not require
elaboration. Kuhn's paradigms are what communities share. The programme has one
author, a handful of student contributors, and an open chat server. A position
held by one person and explained to others is a position. It may be correct. It
is not a disciplinary matrix, because there is no discipline whose matrix it is.
No amount of formal development substitutes for this, and the criterion is not
one a programme can satisfy by wanting to.

\textbf{The anomaly cannot produce a crisis.} This is the more interesting
failure, and it follows from the strongest card in the case for.

Kuhn's crises are generated internally. A paradigm certifies certain puzzles as
solvable, practitioners attempt them using the paradigm's instruments, and the
attempts fail repeatedly in ways the practitioners can see. Representation
dependence does not have this structure. A model that returns different verdicts
for two descriptions of the same case is scored on a benchmark that contains one
of the two descriptions. On that benchmark it is correct. Aggregate accuracy is
insensitive to the defect by construction, because accuracy is computed against
a fixed encoding of each case and never presents the alternative encoding. The
reigning instruments therefore do not register the anomaly as a failure. They do
not register it at all.

This is a real defect in current practice, and identifying it is worth doing.
But it is the wrong shape for a Kuhnian crisis. Kuhn needs a paradigm that is
failing by its own lights, visibly, to the people using it. What is described
here is a paradigm succeeding by its own lights at something its critic regards
as the wrong target. Kuhn has no category for that, and the reason he has none is
that such a complaint cannot build pressure inside the community by the mechanism
his account specifies. It has to be argued from outside, in a vocabulary the
reigning measurements do not contain. Which is what the programme is in fact
doing, and which is not what revolutions look like in \emph{Structure}.

\textbf{The translation procedure is evidence against revolution.} PE-BRW-1.0
specifies how to convert a corpus written under ordinary scholarly norms into a
claim ledger \citep{pebrw}. It defines a characterization step that assigns each
existing claim the class its actual support warrants, a debt ratio measuring the
fraction of claims entered without a prior registration, and a discipline that
forbids re-deriving the corpus under the new method, on the ground that priority
is the one property no later rigour reconstructs. It has been run. A book of roughly thirty-five thousand lines, six foundation papers, and one
further paper converted to 234 claim objects at a debt ratio of 0.962, with five findings and six registrable
predictions identified \citep{casebrownfield}.

Kuhn's incommensurability thesis says no such procedure exists between
paradigms. There is no neutral algorithm that carries the claims of one
disciplinary matrix into another with their standing intact. PE-BRW-1.0 is an
algorithm that does precisely this, it is written down, and its output is
inspectable. Two readings are available. Either the corpus and the ledger are not
separated by a revolution, or Kuhn is wrong about incommensurability. The first
is overwhelmingly more likely.

The right description is that Philosophy Engineering is a conservative extension
of ordinary scholarship. It adds constraints on what may be asserted and
requires records that were not previously required. It does not change the
referents of the terms, and a claim converted into the ledger means what it meant
before. The debt ratio near unity is itself the proof: almost every converted
claim retained its content and gained only a statement about the evidence behind
it. Conservative extensions are valuable. They are not revolutions, and a
programme that can translate its predecessor has by that fact shown it is not
incommensurable with it.

\section{Two arms, and what the second one is}

The programme has three arms. The judgment arm governs what an autonomous system
may assert and contains EIP and BIP. The inquiry arm governs what a research
programme may claim and contains the claim ledger. The discovery arm governs
what a programme may claim survives and contains the transformation registry
\citep{pecls,pedsc}. Much of the plausibility of the paradigm claim comes from
treating these as one thing, and they should be separated.

The judgment arm has the shape of a candidate theory change within one subfield.
It supplies generalizations, changes the criteria of adequacy, and has exemplars.
Whether it becomes a paradigm inside AI alignment is an empirical question about
a community, and the answer is currently no by a wide margin and too early to
assess on any other ground. That is the correct thing to say about a programme of
this age, and saying more would be the error this paper exists to avoid.

The inquiry arm is not a candidate and never was. It contains no theory of any
subject matter. It specifies what a record must contain for certain claims about
that record to be checkable. Reform of this kind has a history, and its history
is not a history of paradigms. The controlled trial, the journal, peer review,
the clinical trial registry, and preregistration each changed what could be
claimed and what had to be shown, and none of them is a paradigm in Kuhn's sense.
The literature on psychology's credibility revolution reached the same verdict
about that episode: substantial reform of evidentiary standards rather than a
change of theory, with the sources calling it a revolution using the word in a
sense other than Kuhn's \citep{vazire2018}.

\subsection{The inquiry arm is a set of secondary rules}

Calling it infrastructure says what it is not. A positive identification is
available and it comes from legal philosophy rather than from philosophy of
science.

Hart distinguishes primary rules, which impose obligations on conduct, from
secondary rules, which govern how primary rules are identified, altered, and
applied \citep[ch.~V]{hart1961}. A normative order with primary rules alone is
what he calls pre-legal, and he holds that such an order suffers three defects.
It is marked by \emph{uncertainty}, since nothing settles which purported rules
are valid ones. It is \emph{static}, since rules change only by slow drift and
never by deliberate act. And it is \emph{inefficient}, since disputes about
whether a rule was broken have no authoritative settlement. The remedies are
three kinds of secondary rule: a rule of recognition, rules of change, and rules
of adjudication.

The claim ledger is a set of secondary rules for a research programme. The
correspondence is close enough to be checked clause by clause.

\begin{center}
\begin{tabular}{p{0.19\textwidth} p{0.19\textwidth} p{0.50\textwidth}}
\toprule
\textbf{Hart's defect} & \textbf{Remedy} & \textbf{What the ledger supplies} \\
\midrule
Uncertainty & rule of recognition & Evidence classes fix what counts as a claim of a given standing, and P3 forbids any document from asserting above the class of what it cites. \\
Static character & rules of change & The registration lifecycle, with \texttt{withdrawn} and \texttt{refuted} carrying distinct propagation, so a commitment can be altered by a recorded act rather than by drift. \\
Inefficiency & rules of adjudication & Blast radius obtained by traversing the declared graph, so what a refutation reaches is settled by computation instead of by argument. \\
\bottomrule
\end{tabular}
\end{center}

This identification does explanatory work the word infrastructure does not. It
predicts a specific failure in corpora that have the first kind of rule and lack
the second, and the prediction is confirmed by the conversion recorded in this
programme's own case studies.

The Geometric Ethics corpus arrived at conversion with a fully developed set of
primary rules. It carried epistemic status tags with published definitions, it
stated that retractions are a feature rather than an embarrassment, and it
required that every revision have provenance \citep{casebrownfield}. Those are
obligations on conduct, stated plainly, and the author observed them. What the
corpus lacked was any rule saying what a correction \emph{does}. Three
falsifications had been recorded, each in the section where it was found, and
none had its consequences propagated, because the dependency column never said
which claims rested on the falsified ones.

Hart's account predicts this outcome exactly, and names it inefficiency: an order
with obligations and no rule of adjudication can recognize a violation and cannot
determine its consequences. The orphan correction is therefore not carelessness
by a careful author. It is the structural signature of primary rules operating
without secondary ones, and it is what the conversion found rather than what the
conversion assumed.

One of the four pathologies falls outside this correspondence, and the mismatch
is worth stating rather than smoothing. Hart's defects concern rules that exist
within a system. The file drawer concerns work that was never entered into one,
and a legal order has no counterpart difficulty, since there is no problem in
jurisprudence corresponding to cases that were never filed. Completeness is
therefore an addition to Hart's scheme rather than an instance of it, which is
consistent with it being the property for which §9 finds no prior art.

\subsection{The precedent is accounting, not astronomy}

The closest historical parallel to a change of this kind is double-entry
bookkeeping. It is not a theory of commerce, it explains nothing about trade, and
no one has proposed it as a paradigm. What it did was make a question
mechanically answerable that had previously rested on a merchant's word. Entering
every transaction twice creates a redundancy whose violation is detectable, so
solvency becomes a property that can be checked rather than asserted. Gap-free
identifier assignment performs the same operation on a different question, namely
whether a programme's asserted results are backed by the evidence it holds.

Two cautions attach to the parallel and both cut against overuse. Sombart's
thesis that double-entry enabled modern capitalism is contested, and the
accounting history literature has been actively revising the evidentiary basis
for claims of that shape \citep{sangster2025}. The parallel is offered for the
mechanism and carries no weight about consequences. The second caution is more
useful. Double-entry spread under external accountability pressure, with
Florentine bankers answerable to their guild for the state of their books, rather
than from merchants wishing to know the truth about themselves. This programme
has no external party entitled to demand its accounts. A method whose historical
analogue succeeded because someone outside could compel it is being operated here
by the only party it constrains, and §10 treats that as the limitation it is.

The consequence worth testing is institutional rather than epistemic. What
auditable books changed was not the conduct of individual merchants but the scale
of enterprise available to them, since capital cannot be pooled among strangers
in a venture whose accounts only its owner can read. If the analogy holds, the
first observable effect of a research ledger is on how many people can contribute
to a programme without a trust relationship to its author, and not on how often
the programme is wrong. Prediction P-V in §8 states this in a form that can fail.

Combining the two arms produces a rhetorical effect that neither supports alone.
The judgment arm lends theoretical content, which reform movements usually lack.
The inquiry arm lends scope across all of science, which a subfield theory lacks.
A reader who takes both at once perceives something resembling a general
scientific revolution. Assessed separately, one is an early theory in a single
field and the other is a set of secondary rules. Neither is a paradigm.

\section{What survives: a present-tense answer to Feyerabend}

The negative result leaves a positive claim standing, and it is more defensible
than the one it replaces.

Lakatos held that a research programme consists of a hard core its adherents
protect by decision and a protective belt of auxiliary hypotheses that absorbs
refutations \citep{lakatos1970}. A programme is progressive when successive
versions predict facts not predicted before and some of those predictions are
confirmed, and degenerating when its modifications only accommodate anomalies
already known. The account rescues scientific rationality from Popper's
difficulty that theories are never abandoned on a single failed prediction, and
it does so by moving the unit of appraisal from the theory to the programme.

Feyerabend's objection is that the appraisal arrives too late to be of use. A
degenerating programme may recover, so no rule tells a scientist when to abandon
one, and the verdict becomes clear only in hindsight \citep{feyerabend1975}. The
Stanford Encyclopedia states the point in the same terms, calling the
methodology pragmatically unsatisfactory because its evaluations are exclusively
retrospective and it therefore cannot advise anyone in the present
\citep{seplakatos}. The objection is usually treated as decisive against
Lakatos's use of the account for demarcation.

The objection depends on a fact about the record rather than a fact about
science. The core and the belt are not marked anywhere at the time the
predictions are made. After a refutation, someone reconstructs which commitments
were central and which were auxiliary, and that person is normally a participant
with an interest in the answer. The reconstruction is not dishonest, and it need
not be. It is simply unconstrained, because no document written before the test
fixed the partition. Kuhn's observation about textbook rewriting is the same
phenomenon at community scale: the record is prose, prose does not carry
dependency structure, and what a refutation actually struck cannot be recovered
from it.

A claim ledger changes that fact about the record. Three of its properties bear
directly on the objection \citep{pecls}. Priority requires that the commit
sealing a registration be an ancestor, in the repository history, of the first
commit that adds the result, so a prediction cannot be written after its
outcome is known. Coherence requires that every dependency a claim declares
resolve to another claim, and that no claim be classified above the weakest
support it rests on. Blast radius is then computed by traversing the declared
graph, so the set of claims a refutation reaches is obtained by traversal rather
than by argument.

The consequence for Lakatos is direct. The core and the belt correspond to
position in the dependency graph, with heavily depended-upon claims occupying the
core and claims nothing rests on occupying the belt. Because the edges are
declared before the test and the ordering is fixed by repository ancestry, the
partition is a property of the record at the time of the seal. It is not
reconstructed and it cannot be adjusted after the outcome is known without
leaving a trace. The question Feyerabend said could only be answered in hindsight,
namely what this refutation strikes, acquires an answer available the moment the
refutation is recorded.

Three limits on this claim must be stated with it.

First, it addresses one input to Lakatos's appraisal rather than the appraisal
itself. Whether a modification is ad hoc still requires judgment about whether a
prediction counted as new, and no graph settles that. The discovery arm's
requirement that transformation families be declared in the registry before the
seal constrains the judgment, since a family cannot be selected after seeing
which tests passed \citep{pedsc}. It does not eliminate it.

Second, localization reaches only as far as the declared conjunction. Duhem and
Quine established that no hypothesis is tested alone, and a refutation indicts
whatever was assumed alongside it \citep{duhem1954,quine1951}. A ledger can
propagate along edges that were written down. Dependencies nobody declared are
invisible to traversal, and the specification says so rather than claiming
otherwise.

Third, the mechanism assumes the repository history is intact. Priority proved
by ancestry is stronger than a third party's timestamp in requiring no trusted
service, and weaker in being vulnerable to history rewriting. The specification
records this trade rather than resolving it, and anchoring a hash in an external
timestamping service closes the gap cheaply.

\section{Predictions}

A paper arguing that records govern appraisal should say what would show it
wrong. Five predictions follow from the accounts in \S6 and \S7, stated so that a
programme other than this one can test them.

\textbf{P-I. Blast radii are small and right-skewed.} If Lakatos is right that
refutations ordinarily strike the belt, and dependencies are declared in advance,
then the number of claims reached by a recorded refutation should be small in
most cases with a long tail. Routinely large radii would indicate either that
the core and belt picture is wrong or that dependencies are under-declared. The
two can be distinguished: under-declaration concentrates large radii on a few
claims with high in-degree, while a wrong picture spreads them evenly.

\textbf{P-II. Recorded negative results rise, then settle high.} A programme
adopting gap-free identifier assignment should record more null and abandoned
lines than matched programmes without a ledger, with the excess concentrated in
the first conversion period as previously undisclosed work receives
dispositions, followed by a steady state above the matched baseline rather than a
return to it.

\textbf{P-III. Wholesale revision loses its discontinuous character.} If the
gestalt quality Kuhn attributes to revolutions is partly an effect of a record
that cannot represent partial change, then programmes with individually
versioned claims should show major revisions as many dated propagations rather
than as a single switch. If ledger programmes still revise discontinuously, the
record was not the cause and Kuhn's psychological account survives intact. This
is the prediction most likely to fail, and it is the one that most directly tests
the paper's thesis.

\textbf{P-IV. Declared partitions diverge from retrospective attribution.} Take
a refutation recorded in a ledger. Ask practitioners who know the field but not
the ledger to say which commitments were central before the test. Compare with
the declaration sealed beforehand. Kuhn's invisibility thesis predicts systematic
divergence, with retrospective attribution making surviving commitments look more
central than they were declared to be. Absence of divergence would indicate that
prose records preserve the partition adequately and that the machinery is
unnecessary.

\textbf{P-V. Scale moves before error rate does.} If the effect of auditable
records is institutional rather than epistemic, as \S6 argues from the accounting
precedent, then the first observable consequence of a programme adopting a ledger
is a change in how many people can contribute to it without a trust relationship
to its author, and not a change in how often its claims turn out wrong. A
programme whose contributor count rises with no corresponding fall in its
retraction rate confirms this. A fall in error rate with no change in
contributor scale refutes it, and would indicate the mechanism is epistemic after
all. The two hypotheses make opposite predictions about the same programme, so a
single case can separate them.

P-IV is inexpensive and can be run on the existing corpus. P-I can be computed
today from any conforming ledger. P-II, P-III, and P-V require a second
programme, and until a programme other than this one adopts the method, none of
the five has been tested.

\section{Related work}

Nothing in \S7 is new in its components, and PRIOR-ART.md states the debts at
length \citep{priorart}. The parts assembled here have the following owners.

Claim-level publication was proposed as nanopublications
\citep{groth2010} and developed as micropublications with argumentation graphs
linking claims to evidence and to each other \citep{clark2014}. Both predate this
specification by a decade or more and both have better serializations.
Manuscripts as version-controlled repositories with continuous integration are
Manubot's contribution \citep{himmelstein2019}. Preregistration established the
norm that priority requires a record made before measurement, and registered
reports made it a condition of publication \citep{chambers2013}. Dependency
resolution and deprecation semantics are borrowed wholesale from package
management, and proof assistant libraries already compute blast radius exactly
for the fragment of claims that can be formalized.

Two adjacent proposals deserve separation. Horstmeyer's Lakat is a
peer-to-peer architecture for continuous integration publishing, explicitly
inspired by Lakatos, with a review-based consensus mechanism and a finality rule
for branched contributions \citep{horstmeyer2023}. Its concern is the governance
of a distributed publishing system, and it does not address the appraisal
question treated in \S7. Vazire's account of psychology's credibility revolution
describes exactly the change in standards of evidence that the inquiry arm
mechanizes, and describes it fifteen years earlier and with a discipline behind
it \citep{vazire2018}. The reading offered here, that such reform is not a
Kuhnian paradigm shift, is not original either. It is the settled view in that
literature.

On infrastructure as a subject for philosophy of science rather than a
precondition of it, the relevant work is Bowker and Star on classification
systems and their invisible consequences \citep{bowker1999}, and Leonelli on how
data infrastructures shape what biology can claim \citep{leonelli2016}. Hacking's
styles of reasoning \citep{hacking1992} and Galison's account of trading zones
and instrument-driven change \citep{galison1997} both describe scientific change
that is not theory change, which is the category this paper places the inquiry
arm in.

Two points about \S6 belong here rather than there. Hart's primary and secondary
rules are standard equipment in legal philosophy and are applied routinely to
international law, to constitutional orders, and to informal normative systems.
I searched for an application to the governance of research programmes and found
none, so \S6 states the correspondence as one I am not aware of prior art for
rather than as one I have established to be new, and a reader who knows of the
application should treat that as a correction to make. The accounting parallel
has been traversed in the opposite direction: at least one study uses Lakatos's
methodology of research programmes to explain why double-entry displaced
single-entry bookkeeping. Borrowing the vehicle back in the other direction is
not an argument, and \S6 uses it only to fix the mechanism.

What is left after these concessions is narrow. The claim is that declaring the
dependency partition before the test, and proving priority by repository
ancestry, converts Lakatos's core and belt from a retrospective reconstruction
into a property of the record, and thereby answers one standing objection to his
methodology. I am not aware of prior art for that specific move. It is also the
only part of this paper that would survive if the rest were withdrawn.

We do not offer a better account of scientific revolutions. We offer the test,
and the test came out negative.

\section{Limits, and what the record cannot show}

The programme's rule is that it may not assert what its ledger cannot show. What
follows is the application of that rule to this paper.

\textbf{One programme, and its author is the subject.} Every observation
supporting \S7 comes from a single corpus, maintained by the person making the
argument, who chose the method and then reported that it worked. The
confirmation risk is severe and cannot be reduced from inside. None of the four
predictions in \S8 has been tested against a programme with different
authorship.

\textbf{The checks do not currently pass.} Running the discovery arm's validator
against the reference programme fails. Ten tests cite three transformation
families the registry never declares. The families were declared in the relevant
sealed registrations and were not mirrored into the registry. The finding is left
unrepaired, because backfilling a registry after the runs is what the declaration
property forbids, and the correct declaration date is not the validator's to
assign \citep{pedsc}. A paper arguing that mechanical checks catch what reading
misses should report the case where they caught the author.

\textbf{Completeness is narrower than it sounds.} Gap-free identifier assignment
makes the set of identified lines of work enumerable. It says nothing about work
that was considered and never given an identifier. The property converts an
unfalsifiable claim into a finite check over a bounded domain, and the boundary
of that domain is set by the same person the property is meant to constrain.

\textbf{The corpus contains falsified predictions that were never propagated.}
Three are recorded: a gauge structure refuted by CHSH-style tests at
\(N=600\), a predicted hysteresis effect that failed to confirm and reversed
under double-blind conditions at \(N=630\), and a predicted schism effect with no
confirmations that failed on the best-powered instrument available. Each was
reported where it was found. None had a blast radius computed, because the
dependency records did not exist when they were found
\citep{casebrownfield}. That is the pathology \S7 describes, observed in the
author's own work before the machinery existed, which is weak evidence that the
machinery was needed and no evidence that it succeeds.

\textbf{No external party can demand these accounts.} The precedent in \S6 cuts
against the programme as well as for it. Double-entry bookkeeping spread where
someone outside the enterprise was entitled to inspect the books and to impose a
cost for their absence. Every institution in the reform history named in \S6 has
the same feature: registries are enforced by journals, trial registration by
regulators, peer review by editors. A ledger operated by the only party it
constrains has no such enforcement, and the properties it checks are therefore
adopted voluntarily by the person they are designed to restrain. Nothing in the
specification fixes this, and it is the most likely way for the method to fail in
practice while appearing to succeed on paper.

\textbf{This paper is an unregistered claim.} It was not sealed before its
argument was developed, it has no instrument, and it has not been verified in
fresh context. Its class is \texttt{exploratory}. Under the assertion authority
property, no document in this corpus may cite it above that class, and a reader
encountering the thesis of \S7 asserted more strongly elsewhere has found a
defect worth reporting.

\section{Conclusion}

The claim tested was that Philosophy Engineering is a new paradigm in Kuhn's
sense. It is not. It has no community, which is the unit Kuhn's account is about.
Its anomaly is invisible to the instruments of the approach it criticizes, so it
cannot generate the internal failure a crisis requires. It possesses a working
procedure for translating its predecessor's corpus into its own vocabulary, and
the existence of that procedure is evidence that the two are not separated by a
revolution at all. The correct description is a conservative extension of
ordinary scholarship, which adds requirements about evidence without changing
what the claims mean.

The claim that survives is smaller. Lakatos's separation of a programme's core
from its protective belt was criticized by Feyerabend for yielding verdicts only
after the fact, and the criticism holds because ordinary scientific records do
not fix the partition in advance. A record that proves priority by ancestry and
requires dependencies to be declared before the test makes the partition a
property of the record, and what a refutation strikes becomes computable when the
refutation is entered rather than arguable afterwards. Section~8 states four ways
that account can fail, and until a programme with different authorship adopts the
method, none of them has been tested.

There is a reason to prefer the smaller claim beyond its being better supported.
Paradigms are announced constantly and the announcements are almost always wrong.
Changes to the record are rare, and the ones that took hold, the journal, the
controlled trial, the registry, changed what could be claimed without anyone
describing them as revolutions. The programme's own rule applies to its
description of itself. It may not assert what its ledger cannot show, and the
ledger does not show a paradigm.

\section*{Acknowledgment of AI use}

Drafting, prior-art search, and structural editing for this paper were performed
with the assistance of a large language model under the author's direction. The
argument, the claims about the corpus, and the verdicts in \S5 are the author's.
Every factual statement about the corpus was checked against the repository
before inclusion. The negative result in \S5 and the limits in \S10 were adopted
on the author's instruction that the paper apply the programme's own assertion
rule to itself.

\begin{thebibliography}{99}
\small

\bibitem[Bowker and Star(1999)]{bowker1999}
Bowker, G.~C. and Star, S.~L.
\newblock \emph{Sorting Things Out: Classification and Its Consequences}.
\newblock MIT Press, Cambridge, MA, 1999.

\bibitem[Bond(2026a)]{foundation}
Bond, A.~H.
\newblock \emph{Philosophy Engineering: A Foundation Document for the
Discipline}, v1.0.
\newblock San Jos\'{e} State University, March 2026.

\bibitem[Bond(2026b)]{pecls}
Bond, A.~H.
\newblock PE-CLS-1.0: The Claim Ledger Specification. Draft, 2026.
\newblock \url{https://github.com/ahb-sjsu/philosophy-engineering}

\bibitem[Bond(2026c)]{pebrw}
Bond, A.~H.
\newblock PE-BRW-1.0: Brownfield Conversion. Draft, 2026.

\bibitem[Bond(2026d)]{pedsc}
Bond, A.~H.
\newblock PE-DSC-1.0: The Discovery Arm. Draft, 2026.

\bibitem[Bond(2026e)]{priorart}
Bond, A.~H.
\newblock Prior art and positioning. Companion to PE-CLS-1.0, 2026.

\bibitem[Bond(2026f)]{casebrownfield}
Bond, A.~H.
\newblock Brownfield conversion of the Geometric Ethics corpus. Case study,
August 2026.

\bibitem[Chambers(2013)]{chambers2013}
Chambers, C.~D.
\newblock Registered reports: A new publishing initiative at \emph{Cortex}.
\newblock \emph{Cortex}, 49(3):609--610, 2013.

\bibitem[Clark et al.(2014)]{clark2014}
Clark, T., Ciccarese, P.~N., and Goble, C.~A.
\newblock Micropublications: A semantic model for claims, evidence, arguments
and annotations in biomedical communications.
\newblock \emph{Journal of Biomedical Semantics}, 5:28, 2014.

\bibitem[Duhem(1954)]{duhem1954}
Duhem, P.
\newblock \emph{The Aim and Structure of Physical Theory}.
\newblock Princeton University Press, 1954. Translation of the 1914 second
edition.

\bibitem[Feyerabend(1975)]{feyerabend1975}
Feyerabend, P.
\newblock \emph{Against Method: Outline of an Anarchistic Theory of Knowledge}.
\newblock New Left Books, London, 1975.
% verify chapter/pages for the Lakatos rhetoric charge on final

\bibitem[Galison(1997)]{galison1997}
Galison, P.
\newblock \emph{Image and Logic: A Material Culture of Microphysics}.
\newblock University of Chicago Press, 1997.

\bibitem[Groth et al.(2010)]{groth2010}
Groth, P., Gibson, A., and Velterop, J.
\newblock The anatomy of a nanopublication.
\newblock \emph{Information Services and Use}, 30(1--2):51--56, 2010.

\bibitem[Hacking(1992)]{hacking1992}
Hacking, I.
\newblock `Style' for historians and philosophers.
\newblock \emph{Studies in History and Philosophy of Science}, 23(1):1--20,
1992.

\bibitem[Hart(1961)]{hart1961}
Hart, H.~L.~A.
\newblock \emph{The Concept of Law}.
\newblock Clarendon Press, Oxford, 1961.
\newblock Primary and secondary rules, and the three defects of a regime of
primary rules, are developed in ch.~V.

\bibitem[Himmelstein et al.(2019)]{himmelstein2019}
Himmelstein, D.~S. et al.
\newblock Open collaborative writing with Manubot.
\newblock \emph{PLOS Computational Biology}, 15(6):e1007128, 2019.

\bibitem[Hohfeld(1913)]{hohfeld1913}
Hohfeld, W.~N.
\newblock Some fundamental legal conceptions as applied in judicial reasoning.
\newblock \emph{Yale Law Journal}, 23(1):16--59, 1913.

\bibitem[Horstmeyer(2023)]{horstmeyer2023}
Horstmeyer, L.
\newblock Lakat: An open and permissionless architecture for continuous
integration academic publishing.
\newblock arXiv:2306.09298, 2023.

\bibitem[Kuhn(1970)]{kuhn1970}
Kuhn, T.~S.
\newblock \emph{The Structure of Scientific Revolutions}, 2nd edition, enlarged.
\newblock University of Chicago Press, 1970.

\bibitem[Kuhn(1974)]{kuhn1974}
Kuhn, T.~S.
\newblock Second thoughts on paradigms.
\newblock In F.~Suppe, editor, \emph{The Structure of Scientific Theories},
pages 459--482. University of Illinois Press, 1974.
% verify pages on final

\bibitem[Lakatos(1970)]{lakatos1970}
Lakatos, I.
\newblock Falsification and the methodology of scientific research programmes.
\newblock In I.~Lakatos and A.~Musgrave, editors, \emph{Criticism and the Growth
of Knowledge}, pages 91--196. Cambridge University Press, 1970.

\bibitem[Leonelli(2016)]{leonelli2016}
Leonelli, S.
\newblock \emph{Data-Centric Biology: A Philosophical Study}.
\newblock University of Chicago Press, 2016.

\bibitem[Masterman(1970)]{masterman1970}
Masterman, M.
\newblock The nature of a paradigm.
\newblock In I.~Lakatos and A.~Musgrave, editors, \emph{Criticism and the Growth
of Knowledge}, pages 59--89. Cambridge University Press, 1970.

\bibitem[Musgrave and Pigden(2025)]{seplakatos}
Musgrave, A. and Pigden, C.
\newblock Imre Lakatos.
\newblock \emph{Stanford Encyclopedia of Philosophy}.
\newblock \url{https://plato.stanford.edu/entries/lakatos/}
% verify edition year and authorship on final

\bibitem[Quine(1951)]{quine1951}
Quine, W.~V.~O.
\newblock Two dogmas of empiricism.
\newblock \emph{The Philosophical Review}, 60(1):20--43, 1951.

\bibitem[Sangster(2025)]{sangster2025}
Sangster, A.
\newblock Time to reboot accounting history: Evidence from scholarship on
double entry.
\newblock \emph{Accounting History}, 2025.
% verify volume, issue and pages on final; cited for the contested status of
% enabling claims about double entry, not for a positive thesis

\bibitem[Vazire(2018)]{vazire2018}
Vazire, S.
\newblock Implications of the credibility revolution for productivity,
creativity, and progress.
\newblock \emph{Perspectives on Psychological Science}, 13(4):411--417, 2018.

\end{thebibliography}

\end{document}
